\documentclass[9pt]{beamer}

\usepackage[T1]{fontenc}
\usepackage[utf8]{inputenc}
\usepackage[frenchb]{babel}
\usepackage{eurosym}

\usepackage{multirow}

\usepackage{graphicx}

\usepackage{ulem}


\newcommand\host{{\it host} }
\newcommand\device{{\it device} }


% --------- BEGIN STYLE BEAMER ---------
\usetheme{Madrid}
\definecolor{ECNBleu}{HTML}{102648}
\definecolor{ECNJaune}{HTML}{FAB600}

\setbeamercolor{frametitle}{bg=ECNBleu,fg=white} %le titre de la frame, en haut
\setbeamercolor{title in head/foot}{bg=ECNJaune,fg=ECNBleu} %barre bas, milieu
\setbeamercolor{author in head/foot}{bg=ECNBleu,fg=ECNJaune} %barre bas, gauche
\setbeamercolor{date in head/foot}{bg=ECNBleu,fg=ECNJaune} %barre bas, droit
\setbeamercolor{section in head/foot}{bg=ECNJaune,fg=ECNBleu} %haut gauche
\setbeamercolor{subsection in head/foot}{bg=ECNBleu,fg=ECNBleu} %haut droite
\setbeamercolor{title}{bg=ECNBleu,fg=white}
\setbeamercolor{block title}{fg=white,bg=ECNBleu}

%\setbeamercolor{normal text}{bg=IUTBleuFonce} %texte
\setbeamercolor{item}{bg=ECNBleu}
\setbeamercolor{subitem}{bg=ECNJaune}
\setbeamercolor{description item}{fg=couleurEmph}
\setbeamercolor{section in toc}{fg=ECNBleu}
%\setbeamercolor{subsubitem}{bg=IUTBleuFonce}

%structure (eq emph dans beamer)
\setbeamercolor{emph}{fg=couleurEmph}
% ----------------------------

\usepackage{listings}

\lstdefinestyle{customc}{
  belowcaptionskip=1\baselineskip,
  breaklines=true,
  frame=L,
  xleftmargin=\parindent,
  language=C,
  showstringspaces=false,
  basicstyle=\tiny\ttfamily,
  keywordstyle=\bfseries\color{green!40!black},
  commentstyle=\itshape\color{purple!40!black},
  identifierstyle=\color{blue},
  stringstyle=\color{orange},
}

\lstset{escapechar=@,style=customc}

\newenvironment<>{questionblock}[1]{%
\setbeamercolor{block title}{bg=ECNJaune,fg=ECNBleu}%
\begin{block}#2{#1}}{\end{block}}

\title[GPGPU]{Programmation sur Processeur Graphique}

\subtitle[\ldots]{
---\\Partie \sout{999} : (Tentative de) Synthèse}
\author[P.-E. Hladik]{Pierre-Emmanuel Hladik\\\tiny pierre-emmanuel.hladik@ec-nantes.fr}
\institute[Centrale Nantes]{}
\date{\today}

\begin{document}

\frame{\titlepage}

%%%%%%%%%%%%%%%%%%%%%%%%%%
\begin{frame}[plain]
\frametitle{Pourquoi utiliser la programmation parallèle ?}

\begin{itemize}
	\item Réaliser un calcul en moins de temps,

	\item Réaliser un calcul dans un laps de temps donné (cas de l'embarqué),

	\item Produire un calcul \og plus précis \fg pour un problème donné dans un laps de temps donné
	
	\item ..
\end{itemize}

\begin{block}{Motivation commune : augmenter la vitesse}
\begin{itemize}
	\item Utiliser plus de ressources de calculs en même temps
	\item Optimiser les performances de calcul = solliciter au maximum les unités de calcul sans temps creux
\end{itemize}
 \end{block}

\end{frame}
%%%%%%%%%%%%%%%%%%%%%%%%%%
\begin{frame}[plain]
\frametitle{... mais pas pour tous les problèmes}

\begin{block}{Quel problème est un bon candidat ?}
\begin{itemize}
	\item Un problème avec beaucoup de données et beaucoup de calculs à chaque itération et/ou beaucoup d'itérations sur les données
	\item Un problème qui puisse être décomposé en sous-problèmes parallèlisables (i.e. les calculs peuvent être réalisés en même temps)
\end{itemize}

 \end{block}



\end{frame}
%%%%%%%%%%%%%%%%%%%%%%%%%%
\begin{frame}[plain]
\frametitle{Le travail du développeur}

\begin{block}{}
Le développeur doit structurer le code et les données de manière à résoudre des sous-problèmes en parallèle.
\end{block}

\begin{block}{Plusieurs approches possibles}

\begin{itemize}
	\item {\color{red}La Bonne} : Accélérer le code initial
	\begin{itemize}
		 \item Utiliser les librairies existantes CuBLAS/CuFFT/thrust/Matlab/PGI pragmas/etc.
	\end{itemize}
	=> travail rapide pour les scientifiques du domaine (connaissances minimales requises)\medskip

	\item {\color{orange}La Mieux} : réécrire/créer du nouveau code
	\begin{itemize}
		\item Permet de repenser l'algorithmique de manière efficace
	\end{itemize}
	=> travail efficace pour les informaticiens (connaissances minimales requises)\medskip

	\item {\color{olive}La Meilleure} : Repenser les méthodes numériques et les algorithmes
	\begin{itemize}
		\item Potentiel avantage en termes de performances
	\end{itemize}
	=> travail interdisciplinaire : nécessite des connaissances en informatique et dans le domaine du problème

\end{itemize}


\end{block}

\end{frame}
%%%%%%%%%%%%%%%%%%%%%%%%%%
%\begin{frame}[plain]
%\frametitle{4 étapes}
%
%décomposition du problème 
%
%choix de l'algorithme
%
%implémentation
%
%Optimisation des performances 
%
%\end{frame}
%%%%%%%%%%%%%%%%%%%%%%%%%%
\begin{frame}[plain]
\frametitle{Comment faire ?}

\begin{block}{Réfléchir, comprendre... puis programmer}
\begin{itemize}
	\item Analyser le problème que vous essayez de résoudre
	\item Comprendre la structure du problème
	\item Appliquer des techniques mathématiques pour trouver la solution	
	\item Adapter le problème à une approche algorithmique (choix de l'algorithme)
	\item Organiser la structure de calcul (décomposition du problème)
	\begin{itemize}
		\item Être conscient de l'interdépendance, des interactions, des goulets d'étranglement
	\end{itemize}
	\item Organiser les données
	\begin{itemize}
		\item Être  conscient de la localité et minimiser les données globales
	\end{itemize}
	\item Enfin, écrivez un code ! (c'est la partie facile... si on oublie le debogage)
\end{itemize}
\end{block}

\end{frame}
%%%%%%%%%%%%%%%%%%%%%%%%%%
\begin{frame}[plain]
\frametitle{Phases types d'un kernel}

\begin{itemize}
	\item Initialisation
	\begin{itemize}
		\item Instancier les structures de données et les peupler
	\end{itemize}
	\item Obtenir un identifiant unique pour chaque thread
	\begin{itemize}
		\item CUDA dispose d'un support intégré pour gérer les index
		\item Calculer un identifiant unique pour chaque thread
		\item Les identifiants sont utilisés pour différencier le comportement des threads et pour les accès mémoire
	\end{itemize}
	
	\item Distribuer les données
	\begin{itemize}
		\item Décomposer les données globales en bloc/tuile et les associer à l'identififant,
		\item Partager et/ou répliquer dans une structure partagée en utilisant l'ID de thread pour associer le sous-ensemble de données aux threads
	\end{itemize}

	\item Exécuter les calculs
	\begin{itemize}
		\item Utiliser les ID des threads pour distribuer les itérations de boucle
		\begin{itemize}
			\item Risque de divergence de la mémoire/des données
		\end{itemize}

		%\item Utiliser l'ID du thread pour réaliser les calculs spécifiques au thread
		
		\item Utiliser l'ID du thread pour faire des actions conditonnées :
		\begin{itemize}
			\item Risque de divergence d'instruction/d'exécution
		\end{itemize}
	\end{itemize}
	\item Reconstruire la structure des données globale, préparer la prochaine itération (si besoin)
\end{itemize}

\end{frame}
%%%%%%%%%%%%%%%%%%%%%%%%%%
\begin{frame}[plain]
\frametitle{Optimisation : focus warp et exécution}

\begin{block}{Gestion des warps}
	\begin{itemize}
		\item Gestion de l'ordonnancement par des warps (groupes de 32 threads) par block :
		\begin{itemize}
			\item Avoir un multiple de 32 (nombre de threads par warp) dans chaque block
			\item Avoir "suffisamment" warps  par bloc pour optimiser l'occupation du SM
		\end{itemize}
		\item Attention aux divergences dans un warp (tous les threads d'un warp exécutent la même instruction)
	\end{itemize}
\end{block}

\begin{block}{Attention aux bouchons}
	\begin{itemize}
		\item Synchronisation des threads
		\item Fonctions Atomiques
	\end{itemize}
\end{block}

\begin{block}{La culture générale c'est important}

Il existe déjà y a des algorithmes dédiés pour traiter certains problèmes, exemple : tuiles, réduction.

\end{block}

\end{frame}
%%%%%%%%%%%%%%%%%%%%%%%%%%
\begin{frame}[plain]
\frametitle{Optimisation : focus sur la gestion de la mémoire}

\begin{block}{Structure classique}
	\begin{itemize}
		\item allocation
		\item copie CPU->GPU
		\item {\it lancement kernel}
		\item copie GPU->CPU
		\item désallocation
	\end{itemize}
\end{block}

\begin{block}{Optimisation gestion mémoire}
	\begin{itemize}
		\item Mémoire partagée : locale à un block, rapide d'accès, utile si réutilisation de la donnée, nécessite une recopie dans cette mémoire avant  le calcul, 
		\item Attention aux accès fusionnés (coalesced)
	\end{itemize}
\end{block}


\begin{block}{Type de gestion de la mémoire}
\begin{itemize}
	\item normale ({\tt malloc} (CPU), {\tt cudaMalloc} (GPU), {\tt cudaMemcpy} (CPU<->GPU)),
	\item pinned ({\tt cudaHostAlloc} (CPU), {\tt cudaMalloc} (GPU), {\tt cudaMemcpy} (CPU<->GPU)),
	\item unifiée ({\tt cudaMallocManaged} (CPU+GPU), {\tt cudaDeviceSynchronize} (CPU))
\end{itemize}
\end{block}

\end{frame}


%%%%%%%%%%%%%%%%%%%%%%%%%%
\begin{frame}[plain]
\frametitle{N'oubliez pas votre trousse à outils}

\begin{itemize}
	\item Vérifier les retours d'erreur (cuda error)
	\item Analyser les performances avec gprof et nvprof
	\item {\color{gray} Debuggeur (gdb)}
\end{itemize}

	\begin{center}
		\includegraphics[width=.5\textwidth]{./figures-pdf/boite_outils}
	\end{center}

\end{frame}



\end{document}